\cotefile{157-3}
\scanpage{34}
\mspage{4}

$W_{\mathbb D}$. Les conditions L1 à L5 sont \unsure{énoncées}
dans \unsure{PS|P5} § 64 (page 164).

\uline{L1} Saturation de $W_{\mathbb D}$ O.K.

\deleted{\unsure{en}} \deleted{\unsure{f}}Cela résulte \inserted{\unsure{formellt}} \deleted{du fait} que l'on \deleted{\unsure{fait}} a
un foncteur
\[
\mathrm{Cat}\xrightarrow{\varphi} C\qquad (=\mathrm{Cat}^{\circ})
\]
tel que \unsure{pour} $u\in\mathrm{Fl}(\mathrm{Cat})$, $u\in W_{\mathbb D}\Longleftrightarrow\unsure{\varphi(u)|\varphi u}$ iso.
(Prendre pour $\varphi(X)$ la \unsure{sous-catégorie} de
$\deleted{\unsure{A}}\mathbb D(X)$ formée des \deleted{\unsure{faisceaux}} \inserted{$\mathbb D$-}coeff. \uline{constants} sur $X$.)

\uline{L2} ``$\mathbb{\Delta}_1$ est $W_{\deleted{\unsure{\mathbb D}}}$-asphérique sur $e$'', i.e. pour tt
$X$ dans $\mathrm{Cat}$, $X\times\mathbb{\Delta}_1\longrightarrow X$ est dans $W$.

Cela signifie aussi que $\mathbb{\Delta}_1\longrightarrow e$ est une $\mathbb D$-fibration
triviale, ce qui est bien le cas.

\uline{L3}, \deleted{(\unsure{renforce} L2) :} Si $X\in\mathrm{Ob}\,\mathrm{Cat}$ a un
objet final, alors $X\longrightarrow e$ est \deleted{$W$-asphérique} \inserted{dans $W$}

% slanted
\margin{En fait $L2\unsure{\Longrightarrow|=}L3$}
\deleted{i.e.} En fait, nous exigeons pour les dérivateurs
que ceci soit vrai \deleted{\unsure{pour}}\inserted{en supposant seulement} $X$ asphérique
(au sens absolu), \deleted{\unsure{seulement},} et on en
conclut \inserted{même} que $X\longrightarrow e$ est $W$-asphérique
i.e. $X\times Y\xrightarrow{\mathrm{pr}_2}Y$ est dans $W$ pour tt $Y$
dans $\mathrm{Cat}$.

\uline{L4} Si $f:X\longrightarrow Y$ est tel que $X_{/y}\longrightarrow Y_{/y}$
sont dans $W$, alors $f$ est dans $W$.
C'est \unsure{OK|OV}.

% archive sheet number: 33
\scanpage{35}
\mspage{5}

\uline{L5} Si $X,Y$ sont \inserted{Cat-}\unsure{fibrées|fibrés} sur $S$, et $X\unsure{\text{---}|\to}Y$ est
un $S$-morphisme \uline{cartésien}, induisant sur les
fibres des $X_s\longrightarrow Y_s$ qui sont dans $W$, alors
$f$ est dans $W$.

\deleted{C'est la} \inserted{Pour $\deleted{\unsure{D}}W_{\mathbb D}$,} prop. 6. page 71 de $\mathrm I$.

Ces axiomes $L_1$ à L5 (équivalents à
L1, L3$'$ ($\mathbb{\Delta}_1\to e$ est dans $W$) et L5) sont dans
la nature d'axiomes de stabilité, assurant
que $W$ est assez grand. Il y a une autre
condition qui limite $W$ à ne pas être trop
grand :

L a) Si $f\in W$, alors $\pi_0(f)$ est bijectif.

Ceci n'est pas satisfait pour tous les $W_{\mathbb D}$
– p.ex. si $D$ est le dérivateur \uline{nul}, alors \deleted{$W_{\mathbb D}$}
$W_{\mathbb D}=\mathrm{Fl}\,\mathrm{Cat}$ ! Mais pour les dérivateurs \unsure{utiles|s'utilisant}
en pratique, cela semble toujours satisfait

L b) c) Le foncteur
\[
\mathrm{Cat}\text{---}\mathrm{Cat}.\,W^{-1}
\]
commute aux sommes \inserted{finies} et aux produits finis$^{*}$.

% slanted
\margin{$^{*}$\unsure{Ces|Les} \unsure{axiomes} si $(\mathrm{Cat},W)$ provient d'un triple de Quillen (avec des Q-fibrations, Q-équiv. faibles, Q-cofibrations)}

À examiner ce que ça signifie en termes
concrets sur $W$ – je n'ai pas vérifié que
c'est OK pour les $W_{\mathbb D}$.

L d) Si $f:X\to Y$ dans $W$, alors $f^\circ:X^\circ\longrightarrow Y^\circ$ aussi,
C'est le cor. 2 page 75 de $(\mathrm I)$,

% archive sheet number: 34
\scanpage{36}

\handwritten{\uline{$\overline{\mathrm{XI}}$} : Hot-fibrations, foncteurs (Hot-)propres et lisses \ldots}

% printed form
% Faint mirrored line-printer text visible through the reused sheet, read in its own direction.
\texttt{\unsure{//*}\quad \unsure{GUTS}\quad 0124\quad I}

\diagram{\begin{tabular}{@{}l@{\quad}l@{}}
\texttt{\phantom{G}GGGGGGGGGG}&\texttt{UU\phantom{UUUUUUUU}UU}\\
\texttt{GGGGGGGGGGGG}&\texttt{UU\phantom{UUUUUUUU}UU}\\
\texttt{GG\phantom{GGGGGGGG}GG}&\texttt{UU\phantom{UUUUUUUU}UU}\\
\texttt{GG}&\texttt{UU\phantom{UUUUUUUU}UU}\\
\texttt{GG}&\texttt{UU\phantom{UUUUUUUU}UU}\\
\texttt{GG}&\texttt{UU\phantom{UUUUUUUU}UU}\\
\texttt{GG\phantom{GGGGG}GGGGG}&\texttt{UU\phantom{UUUUUUUU}UU}\\
\texttt{GG\phantom{GGGGG}GGGGG}&\texttt{UU\phantom{UUUUUUUU}UU}\\
\texttt{GG\phantom{GGGGGGGG}GG}&\texttt{UU\phantom{UUUUUUUU}UU}\\
\texttt{GG\phantom{GGGGGGGG}GG}&\texttt{UU\phantom{UUUUUUUU}UU}\\
\texttt{GGGGGGGGGGGG}&\texttt{UUUUUUUUUUUU}\\
\texttt{\phantom{G}GGGGGGGGGG}&\texttt{\phantom{U}UUUUUUUUUU}
\end{tabular}}

\texttt{*P START STC 124 GUTS\qquad 0070 0001 M06}\\
\texttt{*P START STC 124 GUTS\qquad 0070 0001 M06}\\
\texttt{*P START STC 124 GUTS\qquad 0070 0001 M06}\\
\texttt{*P START STC 124 GUTS\qquad 0070 0001 M06}\\
\texttt{*P START STC 124 GUTS\qquad 0070 0001 M06}\\
\texttt{*P START STC 124 GUTS\qquad 0070 0001 M06}

\texttt{- DERNIERE MISE A JOUR DE JOURNAL : 22-11-82}\\
\texttt{- DEMONSTRATION DU MATERIEL CORAIL (SEMS) LE JEUDI 25/11/82 AU CNUS\unsure{C}}\\
\texttt{- ATTENTION :\quad SUITE A MISE EN ROUTE NOUVELLE VERSION SYSTEME}\\
\texttt{QUELQUES FICHIERS SERONT A RESTAURER OU A RECATAL\unsure{O}}\\
\texttt{SI VOUS NE RETROUVEZ PAS QUELQUES FICHIERS PRENDRE}\\
\texttt{CONTACT AVEC LA PREPARATION.}

\texttt{\unsure{*****************************************************************}}\\
\texttt{* SUR VOS AGENDAS----> REUNION UTILISATEURS CNUSC LE 10/12/82}\\
\texttt{\unsure{*****************************************************************}}

% archive sheet number: 35
\scanpage{37}

% printed form
\figure{\begin{tikzpicture}[x=1mm,y=1mm]
  \begin{scope}[shift={(0,0)}]
  \fill[gray!45] (-2.7,-2.7) rectangle (2.7,2.7);
  \fill[white] (0,0) circle (2.34);
  \draw[gray!65,line width=.15mm] (-2.5,2.6) -- (-1.9,2.6) -- (-1.2,1.7)
    (-1.8,-1.4) -- (-1.2,-2.45) -- (-.62,-1.48)
    (0,3.15) -- (0,2.0) (1.85,0) -- (3.18,0) (0,-2.52) -- (0,-3.05);
  \fill (0,0) circle (1.43);
  \end{scope}
  \begin{scope}[shift={(139,0)}]
  \fill[gray!45] (-2.7,-2.7) rectangle (2.7,2.7);
  \fill[white] (0,0) circle (2.34);
  \draw[gray!65,line width=.15mm] (-2.5,2.6) -- (-1.9,2.6) -- (-1.2,1.7)
    (-1.8,-1.4) -- (-1.2,-2.45) -- (-.62,-1.48)
    (0,3.15) -- (0,2.0) (1.85,0) -- (3.18,0) (0,-2.52) -- (0,-3.05);
  \fill (0,0) circle (1.43);
  \end{scope}
\end{tikzpicture}}

\texttt{//*\quad GUTS\quad 0124\qquad I}

\diagram{\footnotesize\begin{tabular}{@{}l@{\quad}l@{\quad}l@{\quad}l@{}}
\texttt{\phantom{G}GGGGGGGGGG\phantom{G}}&\texttt{UU\phantom{GGGGGGGG}UU}&\texttt{TTTTTTTTTTTT}&\texttt{\phantom{G}SSSSSSSSSS\phantom{G}}\\
\texttt{GGGGGGGGGGGG}&\texttt{UU\phantom{GGGGGGGG}UU}&\texttt{TTTTTTTTTTTT}&\texttt{SSSSSSSSSSSS}\\
\texttt{GG\phantom{GGGGGGGG}GG}&\texttt{UU\phantom{GGGGGGGG}UU}&\texttt{\phantom{GGGGG}TT\phantom{GGGGG}}&\texttt{SS\phantom{GGGGGGGG}SS}\\
\texttt{GG\phantom{GGGGGGGGGG}}&\texttt{UU\phantom{GGGGGGGG}UU}&\texttt{\phantom{GGGGG}TT\phantom{GGGGG}}&\texttt{SS\phantom{GGGGGGGGGG}}\\
\texttt{GG\phantom{GGGGGGGGGG}}&\texttt{UU\phantom{GGGGGGGG}UU}&\texttt{\phantom{GGGGG}TT\phantom{GGGGG}}&\texttt{SSS\phantom{GGGGGGGGG}}\\
\texttt{GG\phantom{GGGGGGGGGG}}&\texttt{UU\phantom{GGGGGGGG}UU}&\texttt{\phantom{GGGGG}TT\phantom{GGGGG}}&\texttt{\phantom{G}SSSSSSSSS\phantom{GG}}\\
\texttt{GG\phantom{GGGGG}GGGGG}&\texttt{UU\phantom{GGGGGGGG}UU}&\texttt{\phantom{GGGGG}TT\phantom{GGGGG}}&\texttt{\phantom{GG}SSSSSSSSS\phantom{G}}\\
\texttt{GG\phantom{GGGGG}GGGGG}&\texttt{UU\phantom{GGGGGGGG}UU}&\texttt{\phantom{GGGGG}TT\phantom{GGGGG}}&\texttt{\phantom{GGGGGGGGG}SSS}\\
\texttt{GG\phantom{GGGGGGGG}GG}&\texttt{UU\phantom{GGGGGGGG}UU}&\texttt{\phantom{GGGGG}TT\phantom{GGGGG}}&\texttt{\phantom{GGGGGGGGGG}SS}\\
\texttt{GG\phantom{GGGGGGGG}GG}&\texttt{UU\phantom{GGGGGGGG}UU}&\texttt{\phantom{GGGGG}TT\phantom{GGGGG}}&\texttt{SS\phantom{GGGGGGGG}SS}\\
\texttt{GGGGGGGGGGGG}&\texttt{UUUUUUUUUUUU}&\texttt{\phantom{GGGGG}TT\phantom{GGGGG}}&\texttt{SSSSSSSSSSSS}\\
\texttt{\phantom{G}GGGGGGGGGG\phantom{G}}&\texttt{\phantom{G}UUUUUUUUUU\phantom{G}}&\texttt{\phantom{GGGGG}TT\phantom{GGGGG}}&\texttt{\phantom{G}SSSSSSSSSS\phantom{G}}
\end{tabular}}

{\tiny
\texttt{*P START STC\ \ 124 GUTS\quad 0070 0001 M06\quad ROOM\quad 1.55.07 PM 23 NOV 82\ \ R6.PR1\quad SYS JES2 START\ \ P*}\\
\texttt{*P START STC\ \ 124 GUTS\quad 0070 0001 M06\quad ROOM\quad 1.55.07 PM 23 NOV 82\ \ R6.PR1\quad SYS JES2 START\ \ P*}\\
\texttt{*P START STC\ \ 124 GUTS\quad 0070 0001 M06\quad ROOM\quad 1.55.07 PM 23 NOV 82\ \ R6.PR1\quad SYS JES2 START\ \ P*}\\
\texttt{*P START STC\ \ 124 GUTS\quad 0070 0001 M06\quad ROOM\quad 1.55.07 PM 23 NOV 82\ \ R6.PR1\quad SYS JES2 START\ \ P*}\\
\texttt{*P START STC\ \ 124 GUTS\quad 0070 0001 M06\quad ROOM\quad 1.55.07 PM 23 NOV 82\ \ R6.PR1\quad SYS JES2 START\ \ P*}\\
\texttt{*P START STC\ \ 124 GUTS\quad 0070 0001 M06\quad ROOM\quad 1.55.07 PM 23 NOV 82\ \ R6.PR1\quad SYS JES2 START\ \ P*}
}

\texttt{- DERNIERE MISE A JOUR DE JOURNAL : 22-11-82}\\
\texttt{- DEMONSTRATION DU MATERIEL CORAIL (SEMS) LE JEUDI 25/11/82 AU CNUSC}\\
\texttt{- ATTENTION :\quad SUITE A MISE EN ROUTE NOUVELLE VERSION SYSTEME}\\
\texttt{QUELQUES FICHIERS SERONT A RESTAURER OU A RECATALOGUER}\\
\texttt{SI VOUS NE RETROUVEZ PAS QUELQUES FICHIERS PRENDRE}\\
\texttt{CONTACT AVEC LA PREPARATION.}

{\scriptsize\texttt{*********************************************************************}\\
\texttt{* SUR VOS AGENDAS----> REUNION UTILISATEURS CNUSC LE 10/12/82\qquad\ *}\\
\texttt{*********************************************************************}}

% margin note written sideways
\margin{MidiGraf Montpellier 21999}

\figure{\begin{tikzpicture}[x=1mm,y=1mm]
  \begin{scope}[shift={(0,0)}]
  \fill[gray!45] (-1.4,-2.9) rectangle (1.45,-1.55);
  \draw[gray!65,line width=.14mm] (1.5,1.5) -- (1.5,-4.5) (0,-2.9) -- (0,-4.5);
  \fill (0,0) circle (1.48);
  \end{scope}
  \begin{scope}[shift={(139,0)}]
  \fill[gray!45] (-1.4,-2.9) rectangle (1.45,-1.55);
  \draw[gray!65,line width=.14mm] (1.5,1.5) -- (1.5,-4.5) (0,-2.9) -- (0,-4.5);
  \fill (0,0) circle (1.48);
  \end{scope}
\end{tikzpicture}}

% archive sheet number: 36
\scanpage{38}
\mspage{\handwritten{100}}

Faut-il croire en l'homéopathie vétérinaire ?

Un dossier de Jacinte Giscard

Sur les 5.000 vétérinaires exerçant en France\unsure{.|,} 4.000
\unsure{utilisent} en certaines occasions des préparations
\unsure{homéopathiques} et 400 d'entre eux traitent 30 \% de
\unsure{leurs} patients uniquement par cette pratique. Les chiffres le révèlent : l'homéopathie vétérinaire est une réalité bien palpable.

Pourtant\unsure{.|,} elle est loin de faire l'unanimité. C'est
même bien souvent un objet de polémique entre partisans qui la portent aux nues et détracteurs qui la font
descendre aux enfers. Parce qu'elle ne s'assoie pas
encore suffisamment sur des preuves scientifiques
indiscutables, l'homéopathie est encore une médecine
qui dérange. Cela est rendu plus vrai du fait qu'elle
s'entoure parfois d'un parfum d'ésotérisme un peu trop
discréditant.

Vous êtes nombreux à nous interroger dans votre
courrier sur l'homéopathie et sur son utilisation chez
l'animal. ``Où en est l'homéopathie vétérinaire, quels
développements peut-on en attendre à l'avenir ?'' sont
des questions qui reviennent chaque jour.

Pour vous, nous avons donc réalisé ce grand dossier, un peu technique parfois (c'était inévitable), mais
riche d'enseignements. Il ne faut pas croire en l'homéopathie comme en une panacée universelle. Mais il est
indubitable qu'elle donne, en de nombreuses occasions, de bons résultats. Essayons de comprendre
pourquoi.

\mspage{47}

\handwritten{``30 millions d'amis''}

\handwritten{nov. \unsure{88|85}}

En fait, et il faut bien se le
mettre dans la tête, l'homéopathie vétérinaire a un intérêt dans toutes les formes
d'intervention médicale. Elle
peut s'appliquer efficacement tant aux cas d'urgence
comme les diarrhées suraiguës ou les traumatismes,
qu'aux maladies courantes
et à la prévention sanitaire.

% archive sheet number: 37
\scanpage{39}

\deleted{\unsure{Hot|Cat}}
\diagram{\begin{tikzcd}[ampersand replacement=\&,row sep=0pt,column sep=2pt,cells={anchor=west},execute at end picture={\draw ([xshift=-2pt]m11.north west) -- ([xshift=-5pt]m11.north west) -- ([xshift=-5pt]m41.south west) -- ([xshift=-2pt]m41.south west);}]
|[alias=m11]| \text{Hot-équivalences} \& \\
\text{foncteurs asphériques, coasphériques : caractérisation} \& \\
\text{foncteurs propres, lisses} \& \\
|[alias=m41]| \text{foncteurs Hot-fibrants,} \& |[draw,rounded corners]| \text{\inserted{Cat-fibrants, Cat-cofibrants}}\arrow[ul,shorten >=1pt]
\end{tikzcd}}

\deleted{foncteurs} fibrations parfaites : $\forall\,\unsure{Y'|y'|y|Y}:\mathbb{\Delta}^1\longrightarrow Y$ \unsure{mor|mon}
\[
\begin{aligned}
X'=X\times_{\unsure{Y|y}}\mathbb{\Delta}^1\longrightarrow\mathbb{\Delta}^1\quad &\text{satisfait :}\quad X'_1\longrightarrow X'\ \text{coasph},\\
&X'_0\longrightarrow X'\ \text{asph.}\quad\text{i.e. les cat.}\\
&x'_0\backslash X'_1\quad(x'_0\in\mathrm{Ob}\,X'_0)\quad\text{et}\quad X'_0/x'_1\quad(x'_1\in X'_1)\\
&\text{sont asphériques.}
\end{aligned}
\]

% archive sheet number: 39
\scanpage{40}
\mspage{127}

1. \uline{Hot-fibrations} \uline{dans} (Cat)

\deleted{\unsure{Proposition 1}} \deleted{\illegible{1}} \inserted{\uline{Théorème 1} :} Soit $f:X\longrightarrow Y$ dans Cat. Conditions équivalentes :

a) Pour tout $Y''\xrightarrow{p}Y'$ sur $Y$, d'où $X''=X\times_Y Y''\xrightarrow{q}X'=X\times_Y Y'$,
si $p$ est une Hot-équivalence, $q$ l'est aussi.

a$'$) Comme a, avec $Y''=e$, $Y'$ ayant un objet final

\deleted{a$''$) Comme \unsure{a)}, avec $Y''=e$, $Y'$ ayant un objet initial}

% slanted
\margin{garder a$''$) et b$''$), conditions duales de a$'$), b$'$)}

b$'$) Pour \unsure{tout} $y\in\mathrm{Ob}\,Y$, $X_y\to X_{/y}$ est une Hot équivalence,
pour \unsure{toute} $y\to y'$ dans $Y$, $X_{/y}\to X_{/y'}$ est une Hot-équivalence,
enfin ces \deleted{\unsure{deux assertions}} conditions sont aussi vérifiées
pour $f':X'\to Y'$, pour tout chgt. de base $Y'\unsure{\text{---}|\to}Y$
dans Cat. (En fait, il suffit de le demander pour la première des
deux conditions, \unsure{nommément|concernant} $X_y\to X_{/y}$.)

\deleted{b$''$) Dual de b$'$, avec $X_y\to {}_{y\backslash}X$ et $ {}_{y'\backslash}X\to {}_{y\backslash}X$}

c) Pour tout faisceau d'ens loc. constant \inserted{$F$} sur $X$,
$f_*(F)$ est loc. constant et se calcule ``fibre par fibre'',
de $=$ \inserted{pour} $R^1f_*(F)$ si $F$ est un faisceau en \unsure{groupes|groupe},
et pour les $\mathrm R^i f_*(F)$ si $F$ est un faisceau abélien. De
plus, on veut que ces conditions restent satisfaites
après tout chgt. de base.

% slanted
\margin{\textcircled{c$'$} version homologique $1^\circ$) en \deleted{\unsure{termes}} d'images \unsure{directes} de groupoïdes $\unsure{f_*|f_!}(F)$, loc. constants, et $2^\circ$) en termes des $\deleted{\unsure{H_i}}\unsure{R^i|R^1} f_!(L)$, pour $L$ faisceau abélien loc. constant.}

\uline{Dém} On a les implications $\pm$ tautologiques
\diagram{\begin{tikzcd}[ampersand replacement=\&,row sep=2pt,column sep=small,execute at end picture={\draw[decorate,decoration={zigzag,segment length=3pt,amplitude=3pt}] (aa.south east) -- (aa2.west) -- (bb2.west);}]
\& a'\arrow[r,Rightarrow]\& b'\arrow[dr,Rightarrow]\&\\
|[alias=aa]| a\arrow[ur,Rightarrow]\arrow[dr,Rightarrow]\&\&\&c,\\
\&|[alias=aa2]| \deleted{\unsure{a''}}\arrow[r,Rightarrow]\&|[alias=bb2]| \deleted{b''}\&
\end{tikzcd}}
et il reste à prouver que $c\Longrightarrow a$.
% slanted
\margin{Comme a est autoduale, étant équivalente à a$'$, b$'$, elle est aussi équivalente aux \uline{conditions duales} de celles-ci a$''$ et b$''$. (cf cor. 2, page 75)}
On est ramené à prouver que si c) est satisfait, et si
\diagram{\begin{tikzcd}[ampersand replacement=\&,row sep=small,column sep=small]
X\arrow[d,"f"']\& X'\arrow[l,"q"']\arrow[d,"f'"]\\
Y\& Y'\arrow[l,"p"]
\end{tikzcd}}
cartésien dans Cat, si $p$ est une Hot-équiv., $q$ aussi.

On utilise le critère d'Artin-Mazur.

$1^\circ$) Si $F$ est un faisceau d'ens. \inserted{loc.} constant sur $Y$,
il faut prouver que $H^0(X,F)\simeq H^0(X',q^*F)$.

% archive sheet number: 40
\scanpage{41}
\mspage{128}

Mais
\[
H^0(X,F)\simeq H^0(Y,\underbrace{f_*(F)}_{\substack{\text{loc. const.}\\\text{par hyp}\\\text{c)}}})
\underset{\substack{\searrow\ \text{car }p\\\text{Hot-équiv.}}}{\xrightarrow{\sim}}
H^0(\deleted{\unsure{X}}Y',p^*\underbrace{f_*(F)}_{\deleted{\unsure{f'_*q^*F}}})\sim\ ,
\]
mais
\[
p^*f_*(F)\xrightarrow{\sim} f'_*(q^*(F)),
\]
car $f_*(F)$ et $f'_*(q^*F)$ se calculent ``fibre par
fibre'' par hypothèse. Donc
\[
H^0(X,F)\simeq H^0(Y',f'_*(q^*(F)))\simeq H^0(X',q^*(F)),\quad\text{\unsure{qed.|cqfd.}}
\]

$2^\circ$) Si $F$ est un faisceau en groupes \inserted{loc.} constant sur $X$,
il faut prouver
\[
H^1(X,F)\xrightarrow{\sim}H^1(X',q^{\unsure{*|\wedge}}F)
\]
Le fait que c'est injectif résulte déjà de $1^\circ$,
qui implique que le foncteur $P\longrightarrow \unsure{f|q}^*P$
\[
\underset{\text{sur }X}{F\text{-torseurs}}\longrightarrow
\underset{\text{sur }X'}{q^*(F)\text{-torseurs}}
\]
est pleinement fidèle, par l'argument bien
connu, \unsure{utilisant} les faisceaux $\mathrm{Hom}$. Il
faut voir que ce foncteur est ess. surjectif.

\deleted{La \deleted{fait que} pleine fidélité implique qu'il
suffit de le voir \uline{localement} sur $Y$, \unsure{d'où} donc
\unsure{OPS|OK si} $Y=Y_{/y_0}$ est muni d'un objet final $y_0$.
Mais \deleted{\illegible{1}}\inserted{\unsure{déjà|alors}} $H^1(X,F)\simeq H^1(X_{/y_0},F)\simeq(R^1 f_*F)_{y_0}\simeq H^0(Y,R^1f_*F)$.
D'ailleurs, par hyp. c) $R^1 f_*(F)$ est loc. constant \deleted{et se
calcule fibre par fibre}, donc comme $p$ une
Hot-équivalence, on a
\[
H^0(Y,R^1 f_*F)\simeq H^0(Y',p^*R^1 f_*F),
\]
}

% archive sheet number: 41
\scanpage{42}
\mspage{129}

\deleted{et comme \unsure{tantôt}, on a
\[
p^*R^1f_*(F)\xrightarrow{\sim}R^1 f'_*(q^*F),
\]
puisque les deux $R^1$ se calculent fibre par fibre.
On trouve donc
\[
H^1(X,F)\simeq H^0(Y',R^1f'_*(q^*F))
\]
\unsure{Or}}

On a \unsure{un hom|une ligne|une liaison} de suites exactes d'ens
ponctués
\diagram{\scriptsize\begin{tikzcd}[ampersand replacement=\&,row sep=small,column sep=small,execute at end picture={\draw ([xshift=-1pt]lasttop.north west) -- ([xshift=-5pt]lasttop.north west) -- ([xshift=-5pt]lastbot.south west) -- ([xshift=-1pt]lastbot.south west);\draw ([xshift=1pt]lasttop.north east) -- ([xshift=5pt]lasttop.north east) -- ([xshift=5pt]lastbot.south east) -- ([xshift=1pt]lastbot.south east);}]
\deleted{\unsure{\leftarrow}}\,H^1(Y,f_*(F))\arrow[r,"u"]\arrow[d,"\alpha"]
\&H^1(X,F)\arrow[r,"v"]\arrow[d,"\beta"]
\&\overset{\unsure{\ni|\in}\,\eta}{H^0(Y,R^1 f_*(F))}\arrow[r,"w"]\arrow[d,"\gamma"]
\&|[alias=lasttop]|H^2(Y,f_*F)\arrow[d,"\delta"']\\
H^1(Y',f'_*(F'))\arrow[r,"u'"]
\&\underset{\unsure{\ni|\downarrow}\,\unsure{\xi'|\zeta'}}{H^1(X',F')}\arrow[r,"v'"]
\&\underset{\downarrow\,\eta'}{H^0(Y',R^1 f'_*(F'))}\arrow[r,"w'"]
\&|[alias=lastbot]|H^2(Y',f'_*F')
\end{tikzcd}}
% slanted
\margin{\unsure{où on a posé} $F'=q^*(F)$}

où les faisceaux $f_*F$, $R^1 f_*(F)$ sont loc constants
par hyp. et se calculent fibre par fibre, les
faisceaux $f'_*(F')$ et $R^1 f'_{\unsure{*|\circ}}(F')$ aussi, ce qui implique
que ces derniers sont can. iso. (par \unsure{les|ces} flèches de
chgt. de base) aux images inverses de leurs
homologues $f_*(F)$, $R^1 f_*(F)$. Comme $p$ est
une Hot-équiv, cela implique que les flèches
$\alpha,\gamma$ sont iso. Donc si $\xi'\in H^1(X',F')$, son image $\eta'$
dans $H^0(Y',R^1 f'_*(F'))$ provient d'un $\eta\in H^0(Y,R^1 f_*(F))$
bien déterminé. Il faut voir que cet $\eta$ provient
d'un $\xi\in H^1(X,F)$. Ceci (si je me rappelle bien)
est possible ssi une certaine obstruction $w(\eta)$
dans $H^2(Y,f_*F)$ est dans le sous-ensemble neutre des
$H^2$ (\uline{NB} \deleted{\unsure{La}} Le champ \inserted{sur $\widehat Y$} des $F$-torseurs sur \inserted{les} $\widehat X/U$
qui remontent $\eta$ est une \deleted{\unsure{ob}} $f_*(F)$ gerbe, \inserted{$\mathscr G_\eta$} et il faut
voir qu'elle est neutre. Mais il y a mieux à faire :

% archive sheet number: 42
\scanpage{43}
\mspage{130}

On \unsure{voit|note} que cette gerbe $\mathscr G$ est loc. constante
(puisque elle correspond à un lien $f_*(\unsure{G|C})$ loc constant),
et que $\mathscr G'$ est son image inverse. Mais comme
$p$ est une Hot-équivalence, il s'ensuit que
le foncteur
\[
\underline{\Gamma}(Y,\mathscr G)\unsure{\text{---}|\to}\underline{\Gamma}(Y',\mathscr G')
\]
est une équivalence de catégories, donc $\unsure{\xi'|\zeta'}$
provient (à iso. près) d'un unique $\xi$ de $\underline{\Gamma}(Y,\mathscr G)$,
qed.

$3^\circ$) Soit $F$ faisceau abélien loc. constant sur $X$,
il faut voir
\[
H^i(X,F)\xrightarrow{\sim} H^i(X',F')\qquad (F'=q^*F),\tag*{$(*)$}
\]
pour tt $i\in\mathbb N$. Mais \deleted{\unsure{comme}} les deux termes sont
l'aboutissement de suites spectrales de
Leray, de termes initiaux
\[
\begin{aligned}
E_2^{pq}&\simeq H^p(Y,R^q f_*(F))\\
E'{}^{pq}_2&\simeq H^p(Y',R^q f'_*(F')
\end{aligned}
\]
et \unsure{ceci|$(*)$} est associé à un hom de suites
spectrales, qui pour les termes initiaux est
\[
H^p(Y,R^q f_*(F))\unsure{\simeq|\longrightarrow}H^{\unsure{q|p}}(Y',R^q f'_*(F'))\tag*{$(**)$}
\]
associé à
\[
p^*(R^q f_*(F))\text{---}R^q f'_*(F')\tag*{$(***)$}
\]
Mais par hyp. les $R^q f_*(F)$, \inserted{$R^q f'_*(F')$}, sont loc const.
et se calculent fibre par fibre,

% archive sheet number: 43
\scanpage{44}
\mspage{131}

d'où il résulte que les hom. de chgt. de base
$(***)$ sont iso, et comme il s'agit de
faisceaux loc. const. et que $p$ est une
Hot-équivalence, il s'ensuit que $(**)$ est
iso, donc aussi $(*)$ sur les aboutissements,
qed.

\uline{Définition 1} Sous les conditions équiv. du th. 1,
on dit que $f$ est une Hot-fibration.

\uline{Proposition 1} Cette notion est stable par chgt. de
base (trivial), par composition (\unsure{formel|formelle} sur a)). Pour
que $f$ soit une Hot-fibration, il f. et s. que $f$
le soit.
\deleted{J'avais envie de mettre les conditions duales
a$''$, b$''$ de a$'$, b$'$, mais il n'y avait pas de place
pour \unsure{elles|aller} dans la démonstration.}

\inserted{\uline{Prop. 2} Soit $f:X\to S$ une Hot-fibration. Alors $s\mapsto\mathrm{hot}(X_s)$ est un système local sur $S$ à val. dans Hot. Donc les $X_s$ pour $s$}
% margin note written sideways
\inserted{\unsure{dans une|d'une} \unsure{$=$} composante connexe de $S$ sont Hot-isomorphes entre eux. (mais pas canoniquement, si cette comp. n'est pas 1-connexe).}

\uline{Théorème 2} \deleted{\unsure{Soit $f$}}
\diagram{\begin{tikzcd}[ampersand replacement=\&,row sep=small,column sep=small]
X^{\unsure{\flat|1}}\arrow[rr,"f"]\arrow[dr,"\unsure{p|p^\flat}"']\&\&X^{\unsure{\sharp|2}}\arrow[dl,"\unsure{p^\sharp|p'}"]\\
\&S\&
\end{tikzcd}}
un diagramme commutatif dans Cat, avec $p,p'$ des Hot-fibrations, et soit $S_0\subset S$ une partie de $S$ telle
que $S_0\to\pi_0(S)$ soit surjectif. Pour que $f$ soit
une \inserted{Hot-}équivalence, il f. et s. que les
$f_s:X_s\unsure{\text{---}|\to}Y_s$ $(s\in S_0)$ le soient.

a) La suffisance se prouve comme tantôt,
avec le critère de Artin-Mazur : \deleted{Il faut}

% archive sheet number: 44
\scanpage{45}
\mspage{132}

Notons que l'hyp sur $S_0$ assure que si $F,F'$ sont
deux faisceaux d'ens. loc. const. sur $S$ et $u:F\unsure{\text{---}|\to}F'$
un \deleted{\unsure{h}}hom \deleted{qui \unsure{sont|est} iso \unsure{sur}} \deleted{\unsure{un}} tel que $u_s:F_s\simeq F'_s$
pour $s\in\unsure{S|S_0}$, alors $u$ est iso.

$1^\circ$) Soit $F$ un faisceau d'ens \inserted{loc} constant sur $X$,
$F'$ son image inverse sur $X'$, il faut prouver
\[
H^0(X,F)\xrightarrow{\sim}H^{\unsure{0|\bullet}}(X',F').
\]
Or
\[
H^0(X,F)\simeq H^0(Y,\underbrace{\unsure{p|f}_*(F)}_{\substack{\text{loc. const}\\\text{sur }S}})\longrightarrow H^0(Y,\unsure{p'|f'}_*(F'))\simeq H^0(X',F'),
\]
il faut prouver que \inserted{$\unsure{(*)|(x)}\quad p_{\unsure{*|=}}(F)\xrightarrow{\sim}p'_*(F')$}
est iso. Or par hyp. sur $p,p'$, les deux
faisceaux sont loc. constants et se calculent
fibre par fibre, d'où résulte que (*) est un
isom pour les fibres en $s\in S_0$, donc un iso
partout. OK

$2^\circ$) Il en résulte déjà, \inserted{si $F$ est un faisceau en \uline{groupes},} que le foncteur
\[
F\text{-torseurs sur }X\text{---}F\text{-torseurs sur }X'
\]
est pl. fidèle, donc $H^1(X,F)\to H^1(X',F')$
injectif, il \deleted{\unsure{f}} reste à prouver que ce foncteur
est ess$^{\mathrm t}$. surjectif. À cause de la pleine
fidélité, c'est \uline{local sur $S$} (\uline{NB} On voit directement que l'hyp. faite sur les $s\in S_0$ implique
la $=$ hyp. pour \uline{tout} $s\in S$, en utilisant le
fait que $p,p'$ \inserted{sont} des fibrations.). Donc on est ramené
au cas où $S=S_{/s_0}$ a un objet final $s_0$.

% slanted
\margin{Les hyps faites sont bel et bien locales sur $S$ (du moins si $S_0=S$)}

% archive sheet number: 45
\scanpage{46}
\mspage{133}

On a alors
\[
\begin{aligned}
H^1(X,F)&\simeq H^0(S,R^1\deleted{\unsure{f}}\unsure{p|f}_*(F))\\
H^1(X',F')&\simeq H^0(S,R^1p'_*(F'))
\end{aligned}
\]
et on voit encore que
\[
R^1p_*(F)\unsure{\text{---}|\to}R^1\deleted{\unsure{f}}p'_*(F')
\]
est iso (car \unsure{ils|il} se \unsure{calculent|calcule} fibre par fibre,
et \deleted{\unsure{alors ils}} on utilise l'hyp. sur les $f_s$),
donc induit un iso
\[
H^1(X,F)\text{---}H^1(X',F').
\]

$3^\circ$) Si $F$ est \deleted{\unsure{un}} abélien, on doit prouver
\[
H^i(X,F)\xrightarrow{\sim}H^i(X',F')
\]
pour tout $i$, et cela résulte encore de la
suite spectrale de Leray pour $(p,F)$ et
pour $(p',F')$.

b) Inversement, supposons que $f$ soit une
Hot-équivalence, et prouvons que les $f_s:X'_s\to X_s$
sont des Hot-équivalences. Je vais utiliser le

\uline{Lemme} Soit $S\in\mathrm{Ob}\,\mathrm{Cat}$, \inserted{$S$ \inserted{0-}connexe.} Alors il existe un
\begin{tabular}{|p{.90\linewidth}}
$\bar S$ \deleted{\unsure{asphé}} $\in\mathrm{Ob}\,\mathrm{Cat}$ \deleted{et}, $\bar S$ Hot-asphérique,
et $\bar S\longrightarrow S$ une Hot-fibration,
\end{tabular}

Voyons comment ceci implique le
th. 2 (partie réciproque) -- \unsure{OPS|O.P.S.} $S$ 0-connexe,
et \deleted{\unsure{on}} fait le chgt de base par $\bar S\to S$

% archive sheet number: 46
\scanpage{47}
\mspage{134}

comme dans le lemme :
\diagram{\begin{tikzcd}[ampersand replacement=\&,row sep=3pt,column sep=3pt]
X'\arrow[rr,"f"]\arrow[ddr,"p'"']\&\&X\arrow[ddl,"p"]\&\&\&\\
\&\&\&\overline{X'}\arrow[rr,"\bar f"]\arrow[ddr,"\overline{p'}"']\&\&\bar X\arrow[ddl,"\bar p"]\\
\&S\&\&\&\&\\
\&\&\&\&\bar S\arrow[ulll,"\unsure{\varphi|q}"']\&
\end{tikzcd}}.

Alors $\overline{p'}$ et $\bar p$ sont des Hot-fibrations, et
comme $\varphi$ est une Hot-fibration, le chgt. de
base par $\varphi$ transforme la Hot-équiv. $f$ en
une Hot-équiv. $\bar f$. Considérons alors
$\bar s\in\mathrm{Ob}\,\bar S$, et le diagramme
\diagram{\begin{tikzcd}[ampersand replacement=\&,row sep=small,column sep=small]
\overline{X'}\arrow[r,"\bar f"]\&\bar X\\
\overline{X'}_{\bar s}\arrow[u,hook,"i'"]\arrow[r,"\bar f_{\bar s}"]\&\bar X_{\bar s}\arrow[u,hook,"i"']
\end{tikzcd}}

Comme $\{\bar s\}\longrightarrow\bar S$ est une Hot-équivalence,
la flèche $i$ qui s'en déduit par le chgt.
de base $\bar p:\bar X\to\bar S$ l'est aussi puisque $\bar p$
est une Hot-fibration, et de $=$ pour
$i'$. Ainsi dans le carré \inserted{comm.} précédent, $\bar f$,
$\unsure{i|j}$ et $i'$ sont des Hot-équivalences, donc aussi
$\bar f_{\bar s}$. Mais \deleted{\unsure{comme}} si $s=\varphi(\bar s)$, $\bar f_{\bar s}$ s'identifie
à $\deleted{\unsure{\bar f_s}}f_s:X'_s\longrightarrow X_s$. D'ailleurs, il est
immédiat que \deleted{\unsure{si $S$}} $\varphi$, étant une Hot-fibration

% archive sheet number: 47
\scanpage{48}
\mspage{135}

à source non vide, et à but connexe, est surjective
(\uline{NB} \unsure{toutes} les fibres d'une $=$ comp. connexe
de $S$ sont Hot-équivalentes\deleted{\unsure{)}}, donc si l'une est non
vide, toutes le sont.) Cela achève la démonstration, modulo le lemme.

Il reste à prouver celui-ci. Comme $S=\varnothing$, je
peux choisir une ``origine''
\[
s_0\in S,
\]
et regarder la catégorie
\[
\underline{\mathrm{Ch}}_{\infty}(\deleted{\unsure{s_0|s}}S;s_0,-)
\]
des chemins infinis dans $S$, d'origine $s_0$, fibre
de $\unsure{\underline{\mathrm{Ch}}|\underline{\mathcal A}}_{\infty}(S)$ par le foncteur-source \deleted{\unsure{$\sigma$}}
\[
\sigma:\underline{\mathrm{Ch}}_{\infty}(S)\longrightarrow S.
\]
On sait que les fibres de $\sigma$ sont Hot-asphériques,
(VII), et la catégorie construite ici n'est
autre \deleted{\unsure{que}} (avec les notations de \uline{$\overline{\mathrm{VII}}$} p.152) que
la \deleted{\unsure{fo}} cat. $\overline{S'}$ qui entre dans la factorisation
canonique de
\[
f:\{s_0\}=S'\unsure{\longrightarrow|\text{---}}S
\]
en
\[
\{s_0\}=S'\xrightarrow[\text{Hot-équiv.}]{}\overline{S'}
 \xrightarrow{\bar f} S,
\]
où ici $\bar f$ s'interprète simplement comme
le foncteur but (i.e le foncteur induit sur

% archive sheet number: 48
\scanpage{49}
\mspage{136}

$\underline{\mathrm{Ch}}_{\infty}(S,s_0,-)$ par le foncteur but
\[
\unsure{\beta|b}:\underline{\mathrm{Ch}}_{\infty}(S)\longrightarrow S\quad).
\]

On a prouvé \deleted{\unsure{dans}} loc. cit. que $\bar f$ est une
Hot-fibration (``propre'' et ``lisse''), mais dans un
contexte qui semblait faire appel tacitement à
la théorie du dérivateur $\mathrm{HOT}$. Donc pour
mettre les résultats ``topologiques'' de loc. cit.
sur une base solide, il faut développer les
notions pertinentes de foncteur propre, foncteur
lisse, sans référence à une théorie $\mathrm{HOT}$. C'est
ce que j'avais d'ailleurs fait déjà \deleted{\unsure{dans}}
en 1983, du temps de \unsure{Pursuing} Stacks,
mais \deleted{sans} rédaction en forme à ce sujet. Il
me faut donc à présent revenir sur ces
notions, de façon un peu systématique. D'ailleurs,
j'aurais du mal à construire un seul foncteur
\deleted{\unsure{propre et lisse}}\inserted{Hot-fibrant} (\unsure{s.f.|sauf} \deleted{\unsure{en}} les exemples triviaux,
tels les isomorphismes\deleted{\unsure{)}}, ou les produits $Y\times Z\xrightarrow{\mathrm{pr}_1}Y$),
en disposant d'une théorie des
foncteurs lisses et des foncteurs propres.$^{*}$
% slanted
\margin{$^{*}$ J'exagère un peu, quand même -- après tout il doit être \unsure{assez évident|clair} que les foncteurs \unsure{qui sont des fibrés au sens de|ou les fibrés au sens de|fibrés au sens de} Kan sont des \deleted{\unsure{Hot-}}\unsure{Hot-fibrés|Hot-fibrants|Hot-fibrations} \unsure{pour le faire|sans faire|pas des tas} (mais il y en a quand même peu!).}

% archive sheet number: 49
